% !PRINTED 251
\weberpage{251}
%<<CONT
und daraus folgt mit Rücksicht auf
\[
\gamma_3(\omega+1) = -\gamma_3(\omega), \quad
\gamma_3\left(-\frac{1}{\omega}\right) = -\gamma_3(\omega): \qquad \text{[§\,54, (14)]}
\]

3. \gesperrt{Ist $n$ ungerade, $c$ durch 4 teilbar, so sind die
Größen}
\[
P^{6}_{c,\eth,a}\,\gamma_3(\omega)^{\frac{n-1}{2}}
\]
\gesperrt{Wurzeln einer Multiplikatorgleichung.} Ist $n \equiv 1 \pmod{4}$, so gilt dasselbe von $P^{6}_{c,\eth,a}$.

Als Beispiel wählen wir $n=3$ und setzen

\weq{5}{x = 27\left(\frac{\eta(3\omega)}{\eta(\omega)}\right)^6, \quad
x_0 = -\left(\frac{\eta\left(\dfrac{\omega}{3}\right)}{\eta(\omega)}\right)^6, \quad
x_1 = -\left(\frac{\eta\left(\dfrac{4+\omega}{3}\right)}{\eta(\omega)}\right)^6, \quad
x_2 = -\left(\frac{\eta\left(\dfrac{8+\omega}{3}\right)}{\eta(\omega)}\right)^6.}

Die Koeffizienten der Gleichung, die diese Größen zu Wurzeln hat,
sind rationale Funktionen von $\gamma_3(\omega)$, und da keine der
Größen (5) für einen endlichen Wert von $\gamma_3$ unendlich wird,
so sind es \gesperrt{ganze} Funktionen von $\gamma_3$. Nach 3. können wir
diese Gleichung also in der Form ansetzen:
\[
x^4 + A_1\gamma_3 x^3 + A_2 x^2 + A_3 \gamma_3 x + A_4 = 0,
\]
worin $A_1, A_2, A_3, A_4$ ganze rationale Funktionen von $j(\omega)$
sind. Zunächst erhält man $A_4$ als das Produkt $x x_0 x_1 x_2$,
welches für keinen Wert von $j(\omega)$ verschwinden kann und daher
konstant sein muß. Aus den Anfängen der Entwickelung [§\,24]

\weq{6}{\begin{aligned}
x &= 27\,q\dots,\\
x_0 &= -q^{-\frac{1}{3}}\dots,\\
x_1 &= -e^{\frac{2\pi i}{3}}q^{-\frac{1}{3}}\dots,\\
x_2 &= -e^{-\frac{2\pi i}{3}}q^{-\frac{1}{3}}\dots,\\
\gamma_3 &= q^{-1}\dots
\end{aligned}}

findet man daher
\[
A_4 = -27.
\]
Es ist ferner
\[
-A_1\gamma_3 = x + x_0 + x_1 + x_2.
\]
Da die rechte Seite für ein unendliches $\gamma_3$, d.\,h. für $q=0$,
nicht einmal in der ersten Ordnung unendlich wird, so muß $A_1$
verschwinden.
